\section{Motywacja i cel pracy}
\label{sec:motywacja}
Stworzenie dokładnego modelu symulacyjnego rakiety sterowanej FOK 2
jest krokiem niezbędnym do powodzenia misji rakiety - dolecenia do celu.
Wymaga to stworzenia modeli ruchu, aerodynamiki, oprogramowania, a także modelu układu wykonawczego systemu sterowania rakiety.
Stworzenie modelu układu wykonawczego wymaga zidentyfikowana modelu, który będzie w stanie odwzorować zachowanie rzeczywistego układu wykonawczego
oraz jego parametrów. Zachowanie całego systemu różni się od zachowania jego poszczególnych elementów, dlatego taki model należy tworzyć w oparciu o badania całego układu.
Ponadto w \cite{olgierd} zidentyfikowano potrzebę wprowadzenia poprawek do każdego 
Została więc zidentyfikowana potrzeba stworzenia stanowiska badawczego, które umożliwi 
powtarzalny pomiar odpowiedzi układu wykonawczego systemu sterowania rakiety w celu zebrania 
danych do identyfikacji modelu.
Jednym z etapów testowania rakiet \cite{cos} jest przeprowadzanie symulacji Hardware-in-the-Loop (HIL).
Ma on na celu weryfikacje modeli symulacyjnych, oprogramowania i integracji systemów. 
Takie testy są tym wierniejsze rzeczywistemu lotowi im bardziej go odzwierciedlają.

O ile w przypadku rakiety FOK 2 możliwe jest wpięcie bezpośrednio komputera pokładowego
rakiety do symulacji, to wymaga to wprowadzenia zmian w konfiguracji oprogramowania względem konfiguracji lotnej.
Może to powodować różnice w zachowaniu systemu a tym samym zmniejszać wierność symulacji.
Dodanie do symulacji HIL rzeczywistego układu wykonawczego systemu sterowania pozwala 
na urzeczywistnienie wyjścia z algorytmu sterowania - wychylenia sterów. Wprowadzenie 
zewnętrznego urządzenia pomiarowego pozwala, natomiast na uniezależnienie się od ograniczeń
wbudowanego w rakietę systemu pomiarowego wychylenia sterów, który jest obarczony błędami 
i ograniczeniami. Ponadto, pozwala to na zmitygowanie opóźnień związanych z przetwarzaniem i transmisją sygnału.

Ostatecznym celem zespołu rozwijającego rakietę FOK 2 w obszarze możliwości przeprowadzania 
symulacji HIL jest stworzenie systemu umożliwiającego na przeprowadzenie misji rakiety lotu do celu
całkowicie w symulacji HIL, bez potrzeby przeprowadzania lotu testowego. Podczas takiej symulacji
rakieta nie powinna \textit{widzieć} różnicy między nią, a rzeczywistym lotem. Do przeprowadzenia 
takiej symulacji do omawianego w tej pracy stanowiska, wymagane jest dołożenie 
stanowiska testowego dla systemu wizyjnego (SPACE) rozwijanego w ramach \cite{Skromak2025HITL}
tworząc tym samym kompletne stanowisko testo HIL rakiety FOK 2 (SHED) \cite{skromak2026hilall}.


\subsection{Cele pracy}
Praca ma stanowić przede wszystkim charakter użytkowy w prowadzeniu badań HIL rakiety FOK 2, a w szczególności jej układu wykonawczego.
Postawiono więc następujące cele pracy:
\begin{itemize}
    \item Stworzenie dokumentacji stanowiska badawczego układu wykonawczego systemu sterowania rakiety FOK 2 (CASKET),
    \item Udokumentowanie symulacji HIL zarówno na komputerze osobistym jak i komputerze czasu rzeczywistego,
    Umożliwienie przeprowadzenia badań HIL rakiety FOK 2 członkom Zespołu.
    \item Identyfikacja modelu układu wykonawczego systemu sterowania rakiety FOK 2.
    \item Przedstawienie ograniczeń, problemów i pól do dalszych ulepszeń stworzonego stanowiska badawczego.
    \item Przedstawienie dalszych kierunków badań HIL rakiety FOK 2.
\end{itemize}